\section{Background}
\label{sec:background}

\subsection{The bittide mechanism and logical synchrony}
bittide realizes logical synchrony through ordered frame exchange
between nodes driven by independent local physical clocks. Logical
synchrony specifies relationships between distributed events without
requiring equal physical phases or a global physical clock~\cite{lall2024logicalSynchrony}.
During normal unstalled operation, nodes send frames to their neighbors
and consume frames from their receiving buffers as their local clocks
advance. Elastic receiving buffers accommodate the difference between
frame delivery and local consumption while preserving frame order.

The bittide controller uses these receiving-buffer occupancies as local
feedback to adjust clock frequency. Faster incoming progress tends to
fill a buffer, whereas faster local consumption tends to drain it;
link-latency changes also alter frame delivery. Frequency control
therefore regulates physical progress while buffer availability supports
the prescribed logical execution. Frequency agreement alone does not
prove the required frame relationships or application-payload correctness.

\subsection{Continuous clocks and buffer dynamics in bittide}

Each node maintains a local physical clock driven by a frequency-controlled oscillator.
In our bittide model, node $i$ has continuous clock phase $\theta_i(t)$ and physical frequency
$\omega_i(t)=\omega_i^u+u_i(t)$, where $\omega_i^u$ is its constant
uncontrolled (uncorrected) clock frequency and $u_i$ is the
applied frequency correction, or control input. In Callisto simulations, each node holds corrections between local actuation events; phase is continuous
and its slope changes at those events. Integer crossings generate
frames inside this continuous plant. Positive physical operation
requires $\omega_i>0$; a negative correction is not a negative frequency.

For edge $j\to i$, occupancy $\beta_{ji}$, initial occupancy
$\beta_{ji}(0)$, arrival count $A_{ji}^{\rm arr}$ and consumption
count $N_i$ obey, while progress guards remain open,
\begin{align}
 \beta_{ji}(t)&=\beta_{ji}(0)+A_{ji}^{\rm arr}(t)-N_i(t),\nonumber\\
 N_i(t)&=\lfloor\theta_i(t)\rfloor-\lfloor\theta_i(0)\rfloor.
 \label{eq:bg_count_balance}
\end{align}
Arrivals include frames already in flight at initialization. Fluid
models replace integer counts by phase increments; at constant physical
link latency $\tau_{ji}$ they give
$\dot{\widetilde\beta}_{ji}=\omega_j(t-\tau_{ji})-\omega_i(t)$.
This analytical approximation is distinguished from actual integer
execution~\cite{lall2022resistance,lall2024bufferOccupancy}.

\subsection{Safety, syntony and buffer centering}
For capacity $C_e$ and correction limit $u_{\max}$, normal operation
seeks $0<\beta_e(t)<C_e$, $\omega_i(t)>0$ and
$|u_i(t)|\le u_{\max}$ along the trajectory. A protection-induced
stall is distinct from successful unstalled control. In the nominal
reference, frequency synchronization (syntony) is expressed by
$\Delta_\omega=\max_i\omega_i-\min_i\omega_i\to0$, together with
occupancy regulation, $\beta_e\to\beta_e^*$, for compatible targets.
At a half-capacity target this is buffer centering; the safety analysis
also varies compatible initial/target occupancy at fixed capacity. These objectives
do not mean phase equality or tracking frequency one. Under persistent
delivery disturbances, we assess residual disagreement and reserves
rather than exact equilibrium. Buffer headroom denotes the remaining
reserve before underflow or overflow.

Existing PI analysis relates performance to gains and resistance
distances~\cite{lall2022resistance}; occupancy and reframing analyses
characterize limiting levels and centering~\cite{lall2024bufferOccupancy,lall2023bufferCentering}.
For bittide, the additional question is how a fully specified nonlinear
local law changes transient buffer use and disturbance response at
controlled scales while respecting the execution constraints above. Appendix~\ref{app:background} expands these distinctions.
